% Part: sets-functions-relations
% Chapter: arithmetization
% Section: cuts
%
\documentclass[../../../include/open-logic-section]{subfiles}

\begin{document}

\olfileid{sfr}{arith}{cuts}
\olsection{Van $\Rat$ naar $\Real$}
In wezen laat de volledigheidseigenschap zien dat elk punt $\alpha$ op
de getallenlijn die lijn precies in twee helften verdeelt: een
onderste helft waarvan $\alpha$ de kleinste bovengrens is en een
bovenste helft waarvan $\alpha$ de grootste ondergrens is. Om de reële
getallen uit de rationale getallen te \emph{construeren}, stelde
Dedekind voor de reële getallen eenvoudig op te vatten als de
\emph{sneden} die de rationale getallen verdelen. We identificeren
$\sqrt{2}$ dus met de \emph{snede} die de rationale getallen
$< \sqrt{2}$ scheidt van de rationale getallen $> \sqrt{2}$.
Laten we dit nauwkeurig formuleren. Als we de rationale getallen in
twee helften verdelen, wordt de gemaakte verdeling al eenduidig bepaald
door haar \emph{onderste} helft. Dit leidt tot de volgende definitie:
\begin{defn}[Snede]
Een \emph{snede} $\alpha$ is een niet-leeg echt beginstuk van de
rationale getallen zonder grootste element. Met andere woorden,
$\alpha$ is een snede dan en slechts dan als:
\begin{enumerate}
	\item \emph{niet-leeg en echt}: $\emptyset \neq \alpha \subsetneq \Rat$
	\item \emph{beginstuk}: voor alle $p,q \in \Rat$ geldt: als $p < q \in \alpha$, dan $p \in \alpha$
	\item \emph{geen maximum}: voor elke $p \in \alpha$ bestaat er een $q \in \alpha$ waarvoor $p < q$
\end{enumerate}
Dan is $\Real$ de verzameling van alle sneden.
\end{defn}
We kunnen nu dus schrijven: $\sqrt{2} = \Setabs{p \in \Rat}{p^2 < 2\text{ of }p < 0}$. Natuurlijk moeten we nog nagaan dat dit inderdaad een
\emph{snede} is, maar dat stellen we uit tot \olref[check]{sec}.
Nu deze objecten zijn gedefinieerd, moeten we, net als eerder, de
elementaire functies en relaties erop definiëren. We beginnen met een
eenvoudig geval:
\begin{align*}
	\alpha \leq \beta \text{ dan en slechts dan als }\alpha \subseteq \beta
\end{align*}
Met deze ordening kunnen we het centrale resultaat \emph{formuleren}:
de verzameling sneden heeft de volledigheidseigenschap. Volledig
uitgeschreven luidt de uitspraak als volgt. Als $S$ een niet-lege
verzameling sneden met een bovengrens is, dan heeft $S$ een kleinste
bovengrens. Preciezer: er is een snede $\lambda$ die een bovengrens van
$S$ is, oftewel\ $(\forall \alpha \in S)\alpha \subseteq
\lambda$, en $\lambda$ is de kleinste van zulke sneden, oftewel\ $(\forall \beta \in \Real)((\forall \alpha \in S)\alpha \subseteq \beta \lif \lambda \subseteq \beta)$. Het
bewijs luidt als volgt:
\begin{thm}\ollabel{realcompleteness}
De verzameling sneden heeft de volledigheidseigenschap.
\end{thm}
\begin{proof}
Zij $S$ een willekeurige niet-lege verzameling sneden met een
bovengrens. Stel $\lambda = \bigcup S$.
%\Setabs{p \in \Rat}{(\exists \alpha \in S)p \in \alpha}$$
We beweren eerst dat $\lambda$ een snede is:
\begin{enumerate}
\item Omdat $S$ niet leeg is, bevat $S$ ten minste één snede,
dus $\emptyset \neq \lambda$. Omdat $S$ een verzameling sneden is,
geldt $\lambda \subseteq \Rat$. Omdat $S$ een bovengrens heeft, is er
een $p \in \Rat$ die in geen enkele snede $\alpha \in S$ voorkomt. Dus
$p\notin \lambda$ en bijgevolg $\lambda \subsetneq \Rat$.
\item Stel dat $p < q \in \lambda$. Dan is er een $\alpha \in S$ met
$q \in \alpha$. Omdat $\alpha$ een snede is, geldt $p \in \alpha$.
Dus $p \in \lambda$.
\item Stel dat $p \in \lambda$. Dan is er een $\alpha \in S$ met
$p \in \alpha$. Omdat $\alpha$ een snede is, is er een $q \in
\alpha$ waarvoor $p < q$. Dus $q \in \lambda$.
\end{enumerate}
Daarmee is de bewering bewezen. Bovendien geldt duidelijk $(\forall \alpha \in S)\alpha
\subseteq \bigcup S = \lambda$, oftewel\ $\lambda$ is een bovengrens van $S$. Stel nu dat $\beta \in \mathbb{R}$ eveneens een bovengrens is, oftewel\ $(\forall \alpha \in S)\alpha \subseteq \beta$. Voor elke $p \in \Rat$ geldt: als $p \in \lambda$, dan is er een $\alpha \in S$ met $p \in \alpha$, zodat $p \in \beta$. Omdat dit voor elk rationaal getal geldt, volgt $\lambda \subseteq \beta$. Dus $\lambda$ is de \emph{kleinste} bovengrens van $S$.
\end{proof}
We hebben dus objecten die aan de volledigheidseigenschap voldoen.
Anders gezegd: onze sneden vertonen geen “gaten”. (Een verdere
“snede” in de reële in plaats van de rationale getallen zou daarom
geen interessante nieuwe objecten opleveren.)
Vervolgens moeten we bewerkingen op de reële getallen definiëren. We
beginnen met een inbedding van de rationale getallen in de reële
getallen door vast te leggen dat $p_\Real =
\Setabs{q \in \Rat}{q < p}$ voor elke $p \in \Rat$. Daarna definiëren
we:
\begin{align*}
	\alpha + \beta &= \Setabs{p + q}{p \in \alpha \land q \in \beta}\\
%	\alpha - \beta &\defis \Setabs{p - q}{p \in \alpha \land q \in \Rat \setminus \beta}\\
	\alpha \times \beta &=
	\Setabs{p \times q}{0 \leq p \in \alpha \land 0 \leq q \in \beta} \cup 0_\Real & \text{als }\alpha, \beta \geq 0_\Real
%	\alpha \div \beta &\defis
%	\Setabs{\nicefrac{p}{q}}{p \in \alpha \land q \in \Rat \setminus \beta}  & \text{if }\alpha \geq 0_\Real, \beta > 0_\Real
\end{align*}
Voor ieder reëel getal leggen we bovendien vast dat nul maal dat getal
en dat getal maal nul beide nul zijn. Voor de overige gevallen van
vermenigvuldiging definiëren we eerst:
\begin{align*}
	-\alpha &= \Setabs{p - q}{p < 0 \land q \notin \alpha}
\end{align*}
en leggen we vervolgens vast:
\begin{align*}
	\alpha \times \beta &\defis
	\begin{cases}
		\mathord{-}\alpha \times \mathord{-}\beta &\text{als }\alpha < 0_\Real\text{ en }\beta < 0_\Real\\
		\mathord{-}(\mathord{-}\alpha \times \beta) &\text{als }\alpha < 0_\Real \text{ en }\beta > 0_\Real\\
		\mathord{-}(\alpha \times \mathord{-}\beta) &\text{als }\alpha > 0_\Real \text{ en }\beta < 0_\Real
	\end{cases}
\end{align*}
Vervolgens moet worden nagegaan dat elk van deze definities steeds een
snede oplevert. Ten slotte moeten we met een eenvoudig (maar omslachtig)
bewijs aantonen dat de aldus gedefinieerde sneden zich precies zo
gedragen als beoogd. Ook dit alles stellen we uit tot
\olref[check]{sec}.
\end{document}